\section{Implications for the Districting Integrity Act}
\label{sec:implications}

\subsection{Fixed Parameters in Statute are Safe}

The central policy conclusion of this paper is that fixing parameters in
the DIA statute is safe for partisan neutrality.
A future legislature could not improve its partisan position by lobbying
for a parameter change, because no parameter change within the studied
range produces a partisan effect exceeding 0.3 seats nationally.

This is a significant finding for the legislative drafting process.
Parameter choices that seemed politically contentious---should $\acounty$
be 3 or 5?---turn out to be inconsequential for partisan outcomes.
The legislature can choose any reasonable value without creating a
partisan controversy.

\subsection{Recommended Statutory Defaults}

Based on the sweep results, we recommend the following statutory defaults:

\begin{table}[h]
\centering
\caption{Recommended DIA statutory default parameters with justification.}
\label{tab:recommended-defaults}
\begin{tabular}{lcc}
\toprule
Parameter & Recommended value & Justification \\
\midrule
$\ufactor$ & 1.0\% & Balances compactness and runtime \\
$\acounty$ & 3.0 & Respects counties without overriding compactness \\
$T$ & 600 & B.16 certified sufficient; runtime $<$60 sec/state \\
$\aswing$ & 1.10 & B.9 regime boundary; best compactness \\
$\wvra$ & 0.4 & Moderate VRA boost; consistent with B.14 \\
\bottomrule
\end{tabular}
\end{table}

\subsection{The Compactness Argument for Small $\ufactor$}

The one parameter with a meaningful compactness effect is $\ufactor$.
Tighter balance ($\ufactor = 0.1\%$) produces $\overline{\PP} = 0.318$
vs.\ $\overline{\PP} = 0.299$ at $\ufactor = 5\%$.
Since the DIA's primary objective is compactness (minimum edge cut), the
statute should prefer smaller $\ufactor$ values.

However, very small $\ufactor$ (e.g., $0.1\%$) can cause METIS to fail
to find any valid partition (insufficient population flexibility), requiring
fallback strategies.
The recommended $\ufactor = 1\%$ is the smallest value that consistently
succeeds across all 50 states.

\subsection{Response to the ``Parameter Manipulation'' Challenge}

Any challenger who argues that the chosen parameters favor one party must
show a parameter combination that:
(1) is within the design space (all five parameters within their statutory ranges),
(2) produces a partisan shift exceeding the baseline seed variance ($\sigma \approx 0.8$ seats nationally from B.7), and
(3) is not already addressed by the statutory default values.

No such combination exists within the ranges studied.
The challenger's only viable argument is that the \emph{algorithm structure}
itself (e.g., the choice of GeoSection over unweighted bisection) is partisan.
This is the argument addressed in B.0, which found that algorithm structure
does affect partisanship---but that the DIA's chosen algorithm (GeoSection
with the default parameters) is as neutral as any algorithmic alternative.

\subsection{Separability of Compactness and Partisan Effects}

A striking feature of the results is the near-complete separability of
compactness sensitivity from partisan sensitivity:
\begin{itemize}
\item $\ufactor$ affects compactness strongly, partisanship weakly.
\item $T$ affects neither compactness nor partisanship (for $T \geq 200$).
\item $\acounty$ affects county integrity strongly, partisanship weakly.
\end{itemize}

This separability is theoretically expected: compactness is a geometric
property that can be optimized within a geographically defined partition,
while partisan outcomes depend on the geographic distribution of partisan
voters, which is not under algorithmic control.
